% Copyright Ben Paul Wise. All Rights Reserved.
% ----------------------------------------------
% Part I -- Overview. Written LAST per the approved drafting order, as a
% management-level summary of Parts II-IV: minimal mathematics, no code,
% every number sourced in Chapter 3.
% ----------------------------------------------
\chapter{Overview}
\label{ch:overview}

\section{What this library is}
\label{sec:o-what}

VINCP is a C++ library for computing solutions of \emph{complementarity
problems} --- a single mathematical shape that captures a remarkable range
of planning and competition questions. In words: some quantities must
balance exactly (conservation laws, accounting identities, probabilities
summing to one); other quantities cannot go negative (shipments,
allocations, prices), and each of those is tied to a matching condition
that must hold whenever the quantity is actually used, with a strict
either/or between ``the quantity is zero'' and ``its condition is tight.''
Optimal plans satisfy this shape (a convex optimum together with its
shadow prices is exactly such a system), and so do competitive equilibria
(each player's ``no profitable deviation'' conditions, side by side). One
problem format therefore serves both a logistics planner and a two-sided
game.

The library's design follows from that unification: one way to pose a
problem, one result format --- the answer, a defect measure (the ``residual'' of
the later chapters), the work done, and an honest converged/not-converged
verdict --- and, between them,
a \emph{menu} of interchangeable solver engines, because no single
numerical method is best (or even viable) across the whole range. Posing,
solving, and judging are kept strictly separate, so an engine can be
swapped without touching the problem, and a new problem can be added
without touching any engine. The full developer's view is
Chapter~\ref{ch:manual}; the mathematics is Chapter~\ref{ch:math}.

\section{Why it exists}
\label{sec:o-why}

The library began as a careful C++ port of working GNU Octave prototypes
--- two classical projection methods and a Newton-type outer loop, with
the reference implementations' numerics preserved and their tolerances
verified. Two escalations drove everything since. First, a production
problem arrived: a distribution-planning optimization over networks of a
few hundred nodes, whose hardest instances simply could not be solved by
the ported methods in any acceptable time. Second, a class of
\emph{game} models arrived from a separate modeling effort --- equilibrium
computations that are provably outside the mathematical guarantees of
every classical method the library had. Each escalation forced new solver
technology; the discipline throughout was that every addition be a
\emph{general-purpose} engine or a reusable composition, never a
one-problem special case, so that the next problem inherits the whole
menu. Chapter~\ref{ch:testing} tells both escalation stories in full.

\section{The problem portfolio}
\label{sec:o-portfolio}

The library is validated against two tiers of problems:

\begin{itemize}
\item \textbf{A flock of minimal problems} --- small instances built
  backward from known answers (random monotone systems, deliberately
  ill-behaved systems, degenerate corner cases). They run in milliseconds,
  check exactly, and pin every engine and safeguard individually; 111
  automated tests stand green at the time of writing.
\item \textbf{Three almost-real-world problems}, each an acceptance
  gate:
  \begin{enumerate}
  \item \emph{Quadratic flow planning} --- a single planner ships goods
    over a network of up to 200 nodes to minimize priority-weighted
    shortfall under supply, delivery, and transport-budget limits; the
    hardest geometry (near-tied route costs) is degenerate by
    construction.
  \item \emph{Strategic allocation of effort} --- a competition in which
    $M$ actors spread limited effort across $N$ options to sway the odds
    of outcomes they value differently; the target is a Nash equilibrium
    --- an allocation from which no actor can gain by unilaterally
    deviating --- verified against an independent industry-standard solver
    (PATH).
  \item \emph{Pre-position and deploy (PPD)} --- a family of two-sided
    military-logistics games (pre-position forces, commit to postures,
    ship under transport budgets, interdict the opponent's shipments) with
    smooth attrition. These models are \emph{deliberately} outside every
    classical solver guarantee, and more versions of increasing
    sophistication are in preparation.
  \end{enumerate}
\end{itemize}

The latter two are games; all three arrive through the same
problem-posing interface, and adding the next one is a documented recipe
(Chapter~\ref{ch:manual}), not a development project.

\section{Solution approach}
\label{sec:o-approach}

Two ideas organize the library.

\paragraph{A layered architecture.} Problem definitions, solver engines,
outer drivers, and engine compositions are separate layers with
one-directional dependencies. The practical payoffs: every engine is
interchangeable behind a common interface; specialized linear algebra can
be plugged into an engine without changing the engine (the decisive move
in the flow-planning story below); and diagnostic instrumentation ---
per-iteration progress, defect breakdowns by named constraint, honest
failure verdicts --- is uniform, so a failed run is a specification for
the next attempt rather than a mystery.

\paragraph{An engine menu matched to problem class.} Classical projection
methods are kept for the problems they are good at (well-conditioned
monotone systems, warm restarts, unusual feasible sets). An
interior-point engine handles the degenerate and large monotone problems
that stall projection methods, and accepts problem-specific fast linear
algebra through a documented seam. A semismooth Newton engine attacks
mixed nonlinear problems directly. And for the game models that defeat
every single engine, the library composes engines: an
\emph{alternating chain} that repeatedly projects back to feasibility,
runs a globally strong engine to its stall, hands off to a locally sharp
engine, remembers the best point ever visited, treats a stage failure as
``switch strategy'' rather than ``abort,'' and --- when progress
deterministically stalls --- restarts from a small, reproducible,
error-proportioned perturbation. The chain solved what nothing else
could, and its every design element traces to a documented failed run
(Chapter~\ref{ch:testing}).

\section{Headline results}
\label{sec:o-results}

\begin{table}[htbp]
\centering
\small
\begin{tabular}{p{0.30\textwidth} p{0.62\textwidth}}
\toprule
\textbf{Problem} & \textbf{Result} \\
\midrule
Flow planning, 200-node hardest geometry, all 14{,}500 routing pairs
  & Solved \textbf{certified-exact in 4.5 seconds} --- the original method
    ran six minutes without an answer, and the first interior-point
    attempt took half an hour for a near-miss; the final configuration is
    $\sim$394$\times$ faster than that attempt \emph{and} $0.4\%$ better
    in objective. \\
\addlinespace
Strategic allocation of effort ($6 \times 10$ instance)
  & The PATH-verified Nash equilibrium reproduced independently by
    \textbf{two different engine classes}, in 83\,ms and 182\,ms. \\
\addlinespace
PPD deployment game (450 unknowns, outside all classical guarantees)
  & Solved to machine precision by the alternating chain in
    $\sim$18--29\,s --- \textbf{twice, finding two distinct verified Nash
    equilibria} --- where every standalone engine fails; those failures
    are kept in the test suite as documentation. \\
\addlinespace
Whole suite
  & \textbf{111 of 111 tests green}, spanning unit, protocol, and
    acceptance layers. \\
\bottomrule
\end{tabular}
\caption{Headline results; sources and full narratives in
Chapter~\ref{ch:testing}.}
\label{tab:headline}
\end{table}

Behind each headline sits a verified reference: the flow-planning answer
carries an optimality certificate by construction; the game answers were
confirmed by the models' author against the GAMS/PATH reference system;
and the equilibrium roster mechanism turns each newly discovered
equilibrium into a one-line, author-verified addition to the acceptance
gate.

\section{Challenges and answers}
\label{sec:o-challenges}

The project's three hard problems, one paragraph each --- with the general
principle each one bought.

\paragraph{Scale and degeneracy (flow planning).} The hardest instances
combine size (tens of thousands of unknowns) with near-tied costs that
make the optimum a wide, flat set rather than a point. Classical
projection methods slow to a crawl precisely there. Screening the problem
down helped the easy geometry and failed the hard one --- and provably
traded away $0.4\%$ of answer quality. The resolution came in two steps:
change the \emph{engine class} (interior-point methods are indifferent to
the flat set --- iteration counts stayed at $\sim$35 regardless of size
and degeneracy), then change the \emph{linear algebra} (a
problem-structured factorization, its algebra verified symbolically in
Maxima before any C++ was written, cut each iteration from $\sim$50
seconds to milliseconds). Principle: match the engine to the problem
class, then make the expensive step exploit the problem's structure ---
and prefer exact-by-construction over approximate-then-certify.

\paragraph{Nonmonotonicity (the games).} The PPD models are designed with
mixed incentives --- committing mass is rewarded, token efforts are wiped
out --- which places them outside every convergence theorem in the
library's arsenal. No single engine survives; each fails in its own
characteristic way. The resolution was \emph{composition}: the engines'
failure modes turn out to be complementary, and a disciplined
alternation --- with feasibility repair between stages, best-point
memory, failures treated as strategy switches, and reproducible
perturbation at deterministic stalls --- reached machine-precision
equilibria twice. Principle: where no theory guarantees a single method,
engineer the composition, keep it bounded and honest, and let its every
element be forced by evidence.

\paragraph{Trust (all three).} Numerical claims are easy to get subtly
wrong. The library's answer is verification in depth: hand-derived
algebra is checked symbolically in Maxima before implementation;
transcribed models are cross-validated by independent engine classes
agreeing on author-verified answers; safeguard logic is tested by
injecting failures rather than hoping to encounter them; and every solver
reports honestly, including the well-defined ``did not converge.''
Principle: an answer is worth what its verification chain is worth.

\section{Status and roadmap}
\label{sec:o-status}

The library stands complete against its current portfolio: all engines
and compositions shipped and documented, all 111 tests green, the three
acceptance problems gated on verified answers. Open items are recorded
with dispositions in Chapter~\ref{ch:testing}; the near-term roadmap, in
expected order:

\begin{itemize}
\item \textbf{More PPD versions} --- larger and tougher models from the
  same modeling effort, each arriving as data plus a builder through the
  documented recipe.
\item \textbf{Comparative benchmark rows} on the hardest flow-planning
  instance, completing this report's performance tables.
\item \textbf{Acceleration wrappers} (Anderson/Halpern) --- designed,
  literature-verified, unbuilt; they would speed every projection engine
  without changing any of them.
\item \textbf{A matrix-free data path} for very large problems, and a
  \textbf{Java (JNI) binding}; both are scoped, neither started.
\end{itemize}
% ----------------------------------------------
% Copyright Ben Paul Wise. All Rights Reserved.
% ----------------------------------------------
